% Compilar com: lualatex slides_arquitetura_cuda.tex (2x)
\documentclass[11pt, aspectratio=169, usepdftitle=false]{beamer}
\usepackage{presentation}
\usepackage[brazil]{babel}
\usepackage{tikz}
\usetikzlibrary{positioning,arrows.meta,calc}
\usepackage{listings}
\newcommand{\vs}{\textit{vs.}\ }
\hypersetup{pdfpagemode=UseNone, hidelinks, pdfdisplaydoctitle=true,
pdftitle={Arquitetura CUDA e overhead do Windows 11 no benchmark de latência}}

\lstset{basicstyle=\ttfamily\scriptsize\setstretch{1}, keywordstyle=\color{NavyBlue}, commentstyle=\color{DarkGray},
  stringstyle=\color{RubineRed}, columns=fullflexible, keepspaces=true, showstringspaces=false,
  backgroundcolor=\color{LightGray!40}, frame=none, xleftmargin=2mm, framexleftmargin=2mm,
  aboveskip=2mm, belowskip=1mm, language=Python, morekeywords={with,as}, literate={á}{{\'a}}1 {é}{{\'e}}1 {ã}{{\~a}}1 {ç}{{\c{c}}}1 {ó}{{\'o}}1 {í}{{\'i}}1 {ú}{{\'u}}1 {õ}{{\~o}}1 {ê}{{\^e}}1}
\tikzset{bx/.style={draw=DarkGray, rounded corners=2pt, fill=LightGray!45, align=center, font=\scriptsize, inner sep=3pt},
  ar/.style={-{Stealth[length=2mm]}, DarkGray, thick}}

\title{Arquitetura CUDA e o overhead do Windows 11}
\author{Benchmark de latência de tokens --- llm-ada-project}
\date{\texttt{timed\_generate} / \texttt{TokenTimer} --- outubro de 2026}

\begin{document}
\frame{\titlepage}

\begin{frame}
\frametitle{Roteiro}
\begin{enumerate}
\item Arquitetura CUDA: abstrações e como a GPU executa (base: CMU 15-779, aulas 3--4)
\item Por que o decode batch=1 é latency-bound e bandwidth-bound
\item Launch, sincronização, streamer, clocks, tensor cores e 4-bit no nosso benchmark
\item \al{Windows 11}: de onde vem o overhead e \al{como calculá-lo}
\item Protocolo de medição e checklist
\end{enumerate}
\vspace{2mm}
{\footnotesize\color{DarkGray} Nada aqui foi executado em GPU: esta máquina não tem CUDA. Códigos e exemplos numéricos são propostas para rodar no Colab T4 / PC Windows.}
\end{frame}

\begin{frame}
\heading{1. Arquitetura CUDA}
\end{frame}

\begin{frame}[fragile]
\frametitle{Abstrações CUDA: grid, bloco, thread}
\begin{columns}[T]
\column{0.5\textwidth}
\begin{itemize}
\item \al{Host} (CPU) lança um \al{kernel}; o \al{device} (GPU) executa
\item Kernel = grid de blocos; bloco = até 1024 threads
\item Cada thread acha seus dados por \texttt{blockIdx}, \texttt{threadIdx}
\item Blocos são \al{independentes}: a GPU os executa em qualquer ordem
\item Lançamento é \al{assíncrono}: a CPU segue em frente
\end{itemize}
\column{0.48\textwidth}
\begin{lstlisting}[language=C, basicstyle=\ttfamily\tiny\setstretch{1}]
__global__ void matrixAdd(float A[Ny][Nx],
    float B[Ny][Nx], float C[Ny][Nx]) {
  int i = blockIdx.x*blockDim.x + threadIdx.x;
  int j = blockIdx.y*blockDim.y + threadIdx.y;
  C[j][i] = A[j][i] + B[j][i];
}
// host
dim3 threadsPerBlock(4, 3);
dim3 numBlocks(Nx/4, Ny/3);
matrixAdd<<<numBlocks, threadsPerBlock>>>(A,B,C);
\end{lstlisting}
{\footnotesize\color{DarkGray} Exemplo adaptado de CMU 15-779 (Z. Jia, 2025).}
\end{columns}
\end{frame}

\begin{frame}
\frametitle{Do bloco ao warp: como a GPU executa}
\begin{columns}[T]
\column{0.52\textwidth}
\begin{itemize}
\item O escalonador mapeia blocos em \al{SMs} enquanto houver recursos (registradores, shared, slots de warp)
\item Bloco é quebrado em \al{warps} de 32 threads (SIMT: mesma instrução)
\item \al{Warp scheduler}: a cada ciclo escolhe um warp pronto; troca de warp esconde a latência da memória
\item \texttt{if} divergente no warp = caminhos serializados
\item T4 (Turing): 40 SMs, 4 schedulers/SM, 32 warps/SM
\end{itemize}
\column{0.46\textwidth}
\begin{tikzpicture}[node distance=3mm]
\node[bx, minimum width=5.6cm] (g) {GPU --- grid de blocos};
\node[bx, below=of g, minimum width=2.6cm, xshift=-1.5cm] (s1) {SM 0\\bloco 0, bloco 2};
\node[bx, below=of g, minimum width=2.6cm, xshift=1.5cm] (s2) {SM 1\\bloco 1, bloco 3};
\node[bx, below=of s1, minimum width=2.6cm] (w) {warps de 32 threads};
\node[bx, below=of s2, minimum width=2.6cm] (ws) {4 warp schedulers};
\node[bx, below=5mm of $(w)!0.5!(ws)$, minimum width=5.6cm] (r) {registradores \;|\; L1/shared 96 KB \;|\; tensor cores};
\draw[ar] (g) -- (s1); \draw[ar] (g) -- (s2); \draw[ar] (s1) -- (w); \draw[ar] (s2) -- (ws);
\end{tikzpicture}
\end{columns}
\end{frame}

\begin{frame}
\frametitle{Hierarquia de memória (CUDA e T4)}
\begin{columns}[T]
\column{0.5\textwidth}
\begin{itemize}
\item Modelo CUDA: privada por thread (registradores), \al{shared} por bloco, \al{global} para todos
\item Shared + \texttt{\_\_syncthreads()}: threads do bloco cooperam (tiling)
\item Em GPUs de datacenter (A100/H100) a global é \al{HBM}; na \al{T4 é GDDR6}, 16 GB, 300 GB/s
\item L2 (4 MB na T4) fica entre os SMs e a DRAM
\end{itemize}
\column{0.46\textwidth}
\centering
\begin{tikzpicture}[node distance=1.5mm]
\node[bx, minimum width=2.4cm] (a) {registradores\\$\sim$ 1 ciclo};
\node[bx, below=of a, minimum width=3.4cm] (b) {L1 / shared (por SM)};
\node[bx, below=of b, minimum width=4.4cm] (c) {L2 4 MB (todos os SMs)};
\node[bx, below=of c, minimum width=5.4cm] (d) {global: GDDR6 300 GB/s (T4) / HBM};
\node[bx, below=of d, minimum width=5.4cm, fill=white] (e) {RAM do host via PCIe (bem mais lenta)};
\end{tikzpicture}
\end{columns}
\end{frame}

\begin{frame}
\frametitle{Tensor cores}
\begin{itemize}
\item Unidade de multiplicação de matrizes dentro do SM (desde Volta; Turing na T4)
\item T4: 65 TFLOPS FP16, 130 TOPS INT8, 260 TOPS INT4 \vs 8,1 TFLOPS FP32 nos CUDA cores
\item Só rendem com matrizes grandes (GEMM): \al{prefill}
\item No \al{decode} batch=1 a conta é matriz $\times$ vetor (GEMV): tensor core fica quase ocioso
\item Da GTX 980 (2014) à H100 (2022): 16 \so 132 SMs e 4,6 \so $\sim$1000 TFLOPS, principalmente por tensor cores (CMU 15-779)
\end{itemize}
\end{frame}

\begin{frame}
\heading{2. O que o benchmark mede}
\end{frame}

\begin{frame}[fragile]
\frametitle{A janela de tempo do \texttt{timed\_generate}}
\begin{lstlisting}
cuda_sync(); start = time.perf_counter()
generated = model.generate(**inputs, ..., use_cache=True, streamer=timer)
cuda_sync(); total = time.perf_counter() - start
ttft = timer.times[0] - start          # 1o token (prefill + 1 passo)
ttlt = timer.times[-1] - start
decode_tokens_per_s = (n_new - 1) / (ttlt - ttft)
\end{lstlisting}
\begin{itemize}
\item Tudo é \al{relógio da CPU}: inclui GPU + Python + driver + SO
\item \texttt{TokenTimer.put}: 1º callback recebe o \al{prompt inteiro} (\texttt{tolist()} + comparação) dentro do TTFT
\item Cada token: \texttt{tokens.cpu()} (cópia D\so H) + callback Python
\item Mesmo código nos 4 scripts (\texttt{qwen35}, \texttt{gemma4}, \texttt{huatuogpt}, \texttt{medgemma})
\end{itemize}
\end{frame}

\begin{frame}
\frametitle{Prefill compute-bound, decode bandwidth-bound}
\begin{columns}[T]
\column{0.55\textwidth}
\begin{itemize}
\item \al{Prefill}: $N$ tokens juntos, GEMM, muita conta por byte lido \so tensor cores
\item \al{Decode}: 1 token por passo, lê \al{todos os pesos} + KV cache para produzir 1 token
\item Piso de tempo por token:
\[ t_{\text{piso}} \approx \frac{\text{bytes de pesos} + \text{bytes de KV}}{\text{banda}} \]
\item Ex.: 4 B parâmetros em 4-bit $\approx$ 2 GB \so $2/300 \approx 6{,}7$ ms/token \so teto $\approx$ 150 tok/s (conta de bolso)
\end{itemize}
\column{0.42\textwidth}
\begin{itemize}
\item Batch=1 e dependência token a token: \al{latency-bound}
\item Poucos blocos por kernel, SMs ociosos, muitos kernels pequenos
\item Por isso o custo \al{fixo por kernel} (launch, driver, SO) aparece no tok/s
\end{itemize}
\end{columns}
\end{frame}

\begin{frame}
\frametitle{Launch overhead, CUDA Graphs e \texttt{torch.compile}}
\begin{itemize}
\item Cada passo de decode = centenas de kernels (camadas $\times$ \{attn, MLP, norm, dequant, LoRA\})
\item Cada launch custa alguns µs de CPU + driver; se o kernel dura poucos µs, a GPU espera a CPU
\item \al{CUDA Graphs}: captura a sequência e relança com uma chamada; exige formas e endereços fixos
\item \texttt{torch.compile(mode="reduce-overhead")} + \texttt{cache\_implementation="static"} usa CUDA Graphs no Transformers
\item No Windows: Inductor precisa de Triton (\texttt{triton-windows}, comunitário) \so validar antes com Unsloth/bitsandbytes/LoRA
\end{itemize}
\end{frame}

\begin{frame}
\frametitle{Sincronização host-device e streamer}
\begin{columns}[T]
\column{0.5\textwidth}
\begin{itemize}
\item \texttt{.cpu()}, \texttt{.item()}, \texttt{.tolist()}, \texttt{bool(tensor)} e \texttt{synchronize()} \al{bloqueiam} a CPU até a GPU terminar
\item O \texttt{generate} já checa parada (\texttt{bool(...)}) a cada token: 1 sync/token com ou sem streamer
\item Streamer adiciona: cópia D\so H + chamada Python por token
\item Cada sync = ida e volta pelo driver: \al{é aí que o Windows pesa}
\end{itemize}
\column{0.46\textwidth}
\begin{tikzpicture}[x=1cm,y=0.55cm]
\node[font=\scriptsize] at (-0.6,1) {CPU}; \node[font=\scriptsize] at (-0.6,0) {GPU};
\foreach \x in {0,2.1,4.2}{
 \fill[NavyBlue!50] (\x,0.8) rectangle +(0.5,0.4);
 \fill[DarkGray!70] (\x+0.5,-0.2) rectangle +(1.1,0.4);
 \fill[RubineRed!60] (\x+1.6,0.8) rectangle +(0.5,0.4);}
\node[font=\tiny, align=left] at (2.6,-1.5) {azul: launch \quad cinza: kernels\\vermelho: sync + streamer (GPU ociosa)};
\end{tikzpicture}
\end{columns}
\end{frame}

\begin{frame}
\frametitle{Clocks (DVFS), warm-up e \texttt{nvidia-smi -lgc}}
\begin{itemize}
\item A GPU sobe e desce o clock conforme carga, temperatura e limite de potência (DVFS)
\item Warm-up atual = \al{8 tokens}: curto para estabilizar clock, alocador e caches de kernel
\item Proposta: 3 gerações completas (160 tokens) descartadas antes de medir
\item \texttt{nvidia-smi -lgc 1590,1590} trava o clock (Linux: root; Windows: admin, depende de GPU/driver; no Colab normalmente sem permissão) \so sempre \al{registrar} clock:
\item[] \texttt{nvidia-smi --query-gpu=clocks.sm,clocks.mem,temperature.gpu,power.draw,} \texttt{clocks\_throttle\_reasons.active --format=csv -lms 200}
\end{itemize}
\end{frame}

\begin{frame}
\frametitle{Quantização 4-bit (NF4) e dequantização}
\begin{itemize}
\item \texttt{LOAD\_IN\_4BIT = True}: pesos em NF4 (bitsandbytes / QLoRA) \so 4$\times$ menos bytes que FP16
\item Decode é bandwidth-bound \so menos bytes ajuda
\item Mas cada matmul precisa \al{dequantizar} para FP16 antes/durante a conta: kernels extras
\item Condição fine-tuned soma \al{LoRA} (r=8) em attn + MLP: mais kernels por token
\item Diferença de tok/s base \vs fine-tuned mistura modelo e \al{número de kernels}: comparar com LoRA mesclado (\texttt{merge\_and\_unload}) ou citar a ressalva
\end{itemize}
\end{frame}

\begin{frame}
\heading{3. Overhead do Windows 11}
\end{frame}

\begin{frame}
\frametitle{Pilha de software: Linux \vs Windows 11}
\centering
\begin{tikzpicture}[node distance=1.5mm]
\node[bx, minimum width=4.6cm] (p) {Python + PyTorch + Transformers};
\node[bx, below=of p, minimum width=4.6cm] (r) {CUDA runtime / driver (modo usuário)};
\node[bx, below=of r, minimum width=4.6cm] (k) {driver NVIDIA no kernel};
\node[bx, below=of k, minimum width=4.6cm] (g) {GPU};
\node[above=1mm of p, font=\small] {Linux};
\node[bx, right=1.8cm of p, minimum width=5.4cm] (p2) {Python + PyTorch + Transformers};
\node[bx, below=of p2, minimum width=5.4cm] (r2) {CUDA runtime / driver (modo usuário)};
\node[bx, below=of r2, minimum width=5.4cm, fill=RubineRed!20] (k2) {\al{WDDM}: dxgkrnl + escalonador de GPU do Windows};
\node[bx, below=of k2, minimum width=5.4cm] (kk2) {driver NVIDIA (KMD)};
\node[bx, below=of kk2, minimum width=5.4cm] (g2) {GPU (compartilhada com o desktop / DWM)};
\node[above=1mm of p2, font=\small] {Windows 11 (WDDM)};
\end{tikzpicture}\par
\vspace{2mm}
{\footnotesize Toda submissão e toda espera passam pela camada extra do WDDM. Em WSL2 há ainda a paravirtualização (VM \so dxgkrnl do host).}
\end{frame}

\begin{frame}
\frametitle{WDDM \vs TCC}
\begin{itemize}
\item Windows tem 2 modelos de driver para GPU NVIDIA (guia de instalação CUDA para Windows):
\begin{itemize}
\item \al{WDDM}: para GPUs com display; é o padrão em GeForce/RTX
\item \al{TCC}: só GPUs sem display (Tesla, alguns Quadro/RTX pro, GTX Titan); usa o modelo WDM
\end{itemize}
\item \alr{GeForce não suporta TCC} (exceto GTX Titan) \so num PC gamer o benchmark sempre roda em WDDM
\item Ver/trocar: \texttt{nvidia-smi -q | findstr "Driver Model"}; \texttt{nvidia-smi -i 0 -dm TCC} (admin, sem monitor, reinicia o driver)
\item TCC tira a GPU do escalonador gráfico: é o mais próximo do Linux dentro do Windows
\end{itemize}
\end{frame}

\begin{frame}
\frametitle{Hardware-Accelerated GPU Scheduling (HAGS) e WSL2}
\begin{itemize}
\item Sem HAGS, uma thread de alta prioridade \al{na CPU} agenda o trabalho da GPU; a Microsoft cita limites de \al{overhead de submissão} e \al{latência até a GPU}
\item Com HAGS, o agendamento vai para um processador dedicado na GPU
\item Ligar/desligar: Configurações \so Sistema \so Tela \so Gráficos \so Configurações padrão \so ``Agendamento de GPU acelerado por hardware'' (reiniciar)
\item NVIDIA (CUDA no WSL2): Submit assíncrono + HAGS reduziram o custo de launch; o gap que resta é em \al{cargas pequenas e sequenciais}, onde a latência de launch não se esconde
\item É exatamente o perfil do \al{decode batch=1} \so o efeito aparece no tok/s, não no prefill
\end{itemize}
\end{frame}

\begin{frame}
\frametitle{Outras fontes de overhead no Windows 11}
\begin{itemize}
\item \al{Sysmem fallback}: com VRAM cheia o driver usa RAM via PCIe sem erro \so tok/s despenca. Painel NVIDIA \so ``CUDA -- Sysmem Fallback Policy'' \so \emph{Prefer No Sysmem Fallback}
\item \al{Desktop na mesma GPU}: DWM, navegador e vídeo dividem a GPU e a VRAM
\item \al{Energia}: plano ``Desempenho máximo''; Painel NVIDIA \so ``Modo de gerenciamento de energia'' \so \emph{Preferir desempenho máximo}; notebook na tomada
\item \al{CPU}: antivírus, indexação, updates disputam o núcleo que lança kernels
\item \al{TDR}: kernel $>$ \texttt{TdrDelay} (padrão 2 s) é abortado --- não muda latência, mas derruba runs longos
\item \al{Relógio} não é o problema: \texttt{time.perf\_counter} usa QPC ($\le$ 1 µs); evite \texttt{time.time()}
\end{itemize}
\end{frame}

\begin{frame}
\frametitle{Como calcular o overhead: decomposição}
Por token de decode $i$:
\[ t_i^{\text{wall}} = \underbrace{t_i^{\text{GPU}}}_{\text{kernels}} + \underbrace{t_i^{\text{launch}} + t_i^{\text{sync}}}_{\text{driver + SO}} + \underbrace{t_i^{\text{py}} + t_i^{\text{streamer}}}_{\text{Python}} \]
\begin{itemize}
\item $t^{\text{wall}}$: \texttt{perf\_counter} (o que o \texttt{TokenTimer} mede)
\item $t^{\text{GPU}}$: CUDA events ou Nsight Systems (relógio da GPU, $\sim$0,5 µs, independe do SO)
\item Overhead de host por token e fração:
\[ O_{\text{host}} = \overline{t^{\text{wall}}} - \overline{t^{\text{GPU}}} \qquad f_{\text{host}} = \frac{O_{\text{host}}}{\overline{t^{\text{wall}}}} \]
\item Streamer: $\Delta_{\text{streamer}} = t^{\text{wall}}_{\text{com}} - t^{\text{wall}}_{\text{sem}}$ (mesmo prompt e tokens)
\end{itemize}
\end{frame}

\begin{frame}
\frametitle{Como calcular o overhead do Windows (atribuição)}
\begin{itemize}
\item Mesma GPU, mesmos pesos, prompt, \texttt{max\_new\_tokens}, versões e clock registrado; mude \al{só o SO}:
\[ \Delta_{\text{Win}} = \tilde t^{\,\text{wall}}_{\text{Windows}} - \tilde t^{\,\text{wall}}_{\text{Linux}} \qquad \%_{\text{Win}} = \frac{\Delta_{\text{Win}}}{\tilde t^{\,\text{wall}}_{\text{Windows}}} \]
($\tilde t$ = mediana por token; Linux = dual boot ou live USB; WSL2 é um terceiro caso, não o baseline)
\item Checagem: se $t^{\text{GPU}}$ é igual nos dois, a diferença é \al{host/driver}, não a GPU
\item Custo por kernel: $\delta_{\text{launch}} = \Delta_{\text{Win}} / K$, com $K$ = kernels por token (contados no Nsight)
\item Ablações para separar causas: HAGS on/off, desktop em outra GPU (iGPU), TCC (se a placa permitir), com/sem streamer
\item Sem Linux disponível: use $O_{\text{host}}$ e a eficiência de banda $\eta = t_{\text{piso}} / \tilde t^{\,\text{wall}}$
\end{itemize}
\end{frame}

\begin{frame}
\frametitle{Exemplo numérico (ilustrativo, não medido)}
\centering
\begin{tabular}{lcc}
\toprule
& Linux & Windows 11 (WDDM) \\
\midrule
$\tilde t^{\,\text{wall}}$ por token & 30,0 ms & 36,0 ms \\
$\tilde t^{\,\text{GPU}}$ por token (events) & 27,0 ms & 27,2 ms \\
$O_{\text{host}}$ & 3,0 ms (10\%) & 8,8 ms (24\%) \\
decode tok/s & 33,3 & 27,8 \\
\bottomrule
\end{tabular}
\begin{itemize}
\item $\Delta_{\text{Win}} = 6{,}0$ ms/token \so $\%_{\text{Win}} = 6/36 \approx 17\%$
\item GPU quase igual (27,0 \vs 27,2) \so a diferença é host/driver/SO
\item Com $K = 600$ kernels/token: $\delta_{\text{launch}} \approx 6\,000/600 = 10$ µs a mais por kernel
\item Se $t^{\text{GPU}}$ também mudar, olhe clock/potência antes de culpar o Windows
\end{itemize}
\end{frame}

\begin{frame}[fragile]
\frametitle{Microbenchmark: latência de launch e de sync}
\begin{lstlisting}
import time, torch
x = torch.zeros(1, device="cuda")

def launch_us(n=10000):  # só enfileira kernels minúsculos
    torch.cuda.synchronize(); t = time.perf_counter()
    for _ in range(n):
        x.add_(1)
    torch.cuda.synchronize()
    return (time.perf_counter() - t) / n * 1e6

def sync_us(n=2000):     # launch + espera, como 1 passo de decode
    torch.cuda.synchronize(); t = time.perf_counter()
    for _ in range(n):
        x.add_(1); torch.cuda.synchronize()
    return (time.perf_counter() - t) / n * 1e6

for _ in range(3):
    launch_us(1000)       # warm-up
print("launch us:", launch_us(), "| launch+sync us:", sync_us())
\end{lstlisting}
{\footnotesize Rodar igual no Windows, WSL2 e Linux. A diferença entre SOs é o custo fixo de driver/SO por kernel e por sync.}
\end{frame}

\begin{frame}[fragile]
\frametitle{Tempo de GPU no benchmark: CUDA events por token}
\begin{lstlisting}
from transformers import LogitsProcessor, LogitsProcessorList

class EventTimer(LogitsProcessor):   # 1 evento por passo, sem sync
    def __init__(self):
        self.events = []
    def __call__(self, input_ids, scores):
        ev = torch.cuda.Event(enable_timing=True); ev.record()
        self.events.append(ev)
        return scores

timer = EventTimer()
start = torch.cuda.Event(enable_timing=True); end = torch.cuda.Event(enable_timing=True)
cuda_sync(); start.record(); t0 = time.perf_counter()
model.generate(**inputs, min_new_tokens=160, max_new_tokens=160, do_sample=False,
               use_cache=True, logits_processor=LogitsProcessorList([timer]))
end.record(); cuda_sync(); wall = time.perf_counter() - t0
gpu = start.elapsed_time(end) / 1000
ms = [start.elapsed_time(ev) for ev in timer.events]
print("wall:", wall, "gpu:", gpu, "O_host:", wall - gpu)
\end{lstlisting}
{\footnotesize Não executado aqui (sem GPU). Events usam o relógio da GPU: lacunas ociosas entre kernels entram no tempo.}
\end{frame}

\begin{frame}
\frametitle{Nsight Systems no Windows: onde a GPU fica parada}
\begin{itemize}
\item \texttt{nsys profile -t cuda,nvtx,osrt -o win python modelos/qwen35.py}
\item \texttt{nsys stats -r cuda\_api\_sum,cuda\_gpu\_kern\_sum win.nsys-rep}
\item $K$ = nº de kernels / tokens gerados; soma das durações = $t^{\text{GPU busy}}$
\item Fração ociosa: $1 - t^{\text{GPU busy}} / t^{\text{wall}}$ \so alta = launch-bound (CPU/driver/SO)
\item \texttt{cuda\_api\_sum}: tempo de \texttt{cudaLaunchKernel} e \texttt{cudaStreamSynchronize} na CPU --- comparar entre Windows e Linux
\item Marcar prefill e decode com \texttt{torch.cuda.nvtx.range\_push/pop}
\end{itemize}
\end{frame}

\begin{frame}
\heading{4. Protocolo e checklist}
\end{frame}

\begin{frame}
\frametitle{Protocolo de medição (Windows 11)}
\begin{enumerate}
\item Registrar: build do Windows, driver, modelo de driver (WDDM/TCC), HAGS, CUDA, PyTorch, Transformers, bitsandbytes, Unsloth, modelo e \texttt{LOAD\_IN\_4BIT}
\item Monitor na iGPU se possível; fechar navegador/jogos; tomada + desempenho máximo; sysmem fallback desligado
\item 1 run frio (anotar à parte) + 3 gerações completas de warm-up descartadas
\item $N \ge 20$ repetições, saída fixa (\texttt{min\_new\_tokens}), \texttt{do\_sample=False}
\item 3 variantes: wall com streamer, wall sem streamer, CUDA events
\item Log de clock/temperatura/potência/throttle em paralelo (\texttt{nvidia-smi -lms 200})
\item Reportar mediana, p90 e IC 95\% (bootstrap); só atribuir ao Windows o que a ablação mostrar
\end{enumerate}
\end{frame}

\begin{frame}
\frametitle{Checklist de otimização do benchmark}
\begin{itemize}
\item \alg{Medir} GPU com CUDA events junto do wall-clock: reportar os dois e $O_{\text{host}}$
\item \alg{Warm-up maior}: 3 $\times$ 160 tokens em vez de 8
\item \alg{Medir sem streamer} e reportar $\Delta_{\text{streamer}}$; prompt grande pesa no TTFT pelo 1º \texttt{put}
\item \alg{Fixar/registrar clocks}; descartar runs com throttle
\item Comparar base \vs fine-tuned com LoRA \al{mesclado} para não medir kernels extras
\item Piloto de \texttt{torch.compile} + cache estático (CUDA Graphs) para cortar launch overhead
\item No Windows: HAGS on, sysmem fallback off, desktop fora da GPU; se puder, Linux ou TCC
\end{itemize}
\end{frame}

\begin{frame}
\frametitle{Onde achar slides oficiais e referências}
\scriptsize\setstretch{1.05}
\begin{itemize}
\item NVIDIA GTC On-Demand, série S. Jones: \url{https://www.nvidia.com/en-us/on-demand/session/gtcfall22-a41101/} (How CUDA Programming Works), \url{https://www.nvidia.com/en-us/on-demand/session/gtcspring23-S51210/}, \url{https://www.nvidia.com/en-us/on-demand/session/gtc24-s62401/}, \url{https://www.nvidia.com/en-us/on-demand/session/gtc25-s72897/}
\item OLCF/NERSC CUDA Training Series: \url{https://www.olcf.ornl.gov/cuda-training-series/}; DLI Teaching Kits: \url{https://developer.nvidia.com/teaching-kits}
\item CMU 15-779 L3--4 (Z. Jia, 2025): \url{https://www.cs.cmu.edu/~zhihaoj2/15-779/slides/03-CUDA-programming.pdf}
\item Turing whitepaper e T4 datasheet (nvidia.com); CUDA Best Practices: \url{https://docs.nvidia.com/cuda/cuda-c-best-practices-guide/}
\item WDDM/TCC: \url{https://docs.nvidia.com/cuda/cuda-installation-guide-microsoft-windows/}; \texttt{nvidia-smi}: \url{https://docs.nvidia.com/deploy/nvidia-smi/}
\item HAGS: \url{https://devblogs.microsoft.com/directx/hardware-accelerated-gpu-scheduling/}
\item CUDA no WSL2: \url{https://docs.nvidia.com/cuda/wsl-user-guide/}, \url{https://developer.nvidia.com/blog/leveling-up-cuda-performance-on-wsl2-with-new-enhancements/}
\item Sysmem fallback: \url{https://nvidia.custhelp.com/app/answers/detail/a_id/5490}; TDR: \url{https://learn.microsoft.com/en-us/windows-hardware/drivers/display/tdr-registry-keys}
\item QPC: \url{https://learn.microsoft.com/en-us/windows/win32/sysinfo/acquiring-high-resolution-time-stamps}; Nsight Systems: \url{https://docs.nvidia.com/nsight-systems/UserGuide/}
\item PyTorch CUDA semantics: \url{https://docs.pytorch.org/docs/stable/notes/cuda.html}; HF LLM optims: \url{https://huggingface.co/docs/transformers/llm_optims}
\end{itemize}
\end{frame}

\lastslide
\end{document}
